\capituloPregunta{¿Cómo convertir datos dispersos en información confiable?}

\section{Motivación}

Los datos no hablan solos. Antes de comparar tratamientos, el estudiante debe entender que toda medición contiene variación. A veces la variación es parte natural del fenómeno; otras veces aparece por errores de medición, duplicación de registros o falta de control.

\section{Caso integrado: gatos y muestreo}

El material del curso contiene una práctica sobre conteo de gatos en la UAM. La pregunta parece sencilla: ¿cuántos gatos hay? Sin embargo, al intentar responderla aparecen problemas reales de investigación:

\begin{itemize}
  \item dos estudiantes pueden contar al mismo gato;
  \item algunos gatos aparecen en más de un lugar;
  \item el horario modifica la probabilidad de observación;
  \item el observador puede confundir individuos similares;
  \item el conteo depende de una regla antiduplicación.
\end{itemize}

\begin{idea}
El muestreo no es una formalidad. Es la forma en que decidimos qué parte de la realidad será observada y bajo qué reglas.
\end{idea}

\section{De observaciones a información}

Una tabla de datos debe permitir distinguir:

\begin{itemize}
  \item patrón central;
  \item dispersión;
  \item valores atípicos;
  \item registros duplicados;
  \item datos faltantes;
  \item condiciones de observación.
\end{itemize}

\section{Actividad}

Diseñe una regla para contar gatos evitando duplicaciones. Después responda:

\begin{enumerate}
  \item ¿Cuál es la unidad de observación?
  \item ¿Qué información mínima se necesita para no contar dos veces?
  \item ¿Qué sesgo aparece si se cuenta solo en un horario?
  \item ¿Qué gráfica ayudaría a explicar la variabilidad?
\end{enumerate}

\section{Conexión}

Cuando ya sabemos observar variabilidad, podemos preguntar si varios tratamientos producen respuestas realmente distintas.
